\section*{Chapter \ref{ch:ml:feature-engineering-selection-and-importance} --- Feature Engineering, Selection and Importance}

\begin{solution}{exo:ml:feature-engineering-selection-and-importance:1}
$10^2/((2\times0.6 - 1)\times200) = 2.5$ expected false selections at 0.6; $100/(0.8\times200) = 0.625$ at 0.9.
\end{solution}

\begin{solution}{exo:ml:feature-engineering-selection-and-importance:2}
Averaging marginal contributions over the six orderings: player~1 contributes $1, 1, 2, 2, 1, 2$, so $\phi_1 = 1.5$;
player~2 likewise $\phi_2 = 1.5$; player~3 adds nothing in any ordering, $\phi_3 = 0$. They sum to $v(P) = 3$.
\end{solution}

\begin{solution}{exo:ml:feature-engineering-selection-and-importance:3}
0.09 points. It does not show that the feature is unimportant: another feature may stand in for it (the \omterm{def:ml:feature-engineering-selection-and-importance:substitution}{substitution
effect}); only the importance of its cluster answers that.
\end{solution}

\begin{solution}{exo:ml:feature-engineering-selection-and-importance:4}
The model uses both copies. Permuted together, both feed it noise with the variance of the signal they carried, so its
predictions become noise of the same size, and a noisy prediction scores worse than predicting the mean: R-squared falls
to $2.67 - 2.72 = -0.05\%$.
\end{solution}

\begin{solution}{exo:ml:feature-engineering-selection-and-importance:5}
MDI adds up the training-loss reduction of every split, and on training data a split on noise does reduce the loss
(it fits noise). \omterm{def:ml:feature-engineering-selection-and-importance:mdi}{Permutation importance} is measured on held-out rows, where those splits help nothing.
\end{solution}

\begin{solution}{exo:ml:feature-engineering-selection-and-importance:6}
One-at-a-time pruning removes correlated copies one after the other: each looks unimportant while the other remains, and
the second removal then costs what neither looked worth (1.05 points here). Prune clusters, not features. If the
importances were computed on the test data, the selection is also a leak (chapter~3).
\end{solution}

\begin{solution}{exo:ml:feature-engineering-selection-and-importance:7}
\texttt{heavy\_tail(df=5.0)}: raw 1.36\%, winsorised 1.49\%, ranked 1.57\%, truth 1.79\%. With five degrees of freedom
extreme values are rare and moderate, so they no longer set the slope; the transforms matter much less than with 1.5.
\end{solution}

\begin{solution}{exo:ml:feature-engineering-selection-and-importance:8}
If the prediction depends on $x_1 + x_{1b}$ only and the two are identical, the value function is symmetric in them, and
the symmetry axiom gives equal \omterm{def:ml:feature-engineering-selection-and-importance:shapley}{Shapley values}. Take $p = \gamma(x_1 + x_{1b}) = 2\gamma x_1$ with $\Var(x_1) = \sigma^2$
and residuals uncorrelated with the features. Permuting $x_1$ alone replaces $p$ by $\gamma(x_{1,\pi} + x_1)$ and adds
$\E[\gamma^2(x_1 - x_{1,\pi})^2] = 2\gamma^2\sigma^2$ to the mean squared error; the same for $x_{1b}$; permuting both with
one permutation adds $\E[4\gamma^2(x_1 - x_{1,\pi})^2] = 8\gamma^2\sigma^2$, twice the sum of the singles.
\end{solution}

\begin{solution}{pb:ml:feature-engineering-selection-and-importance:1}
\begin{enumerate}
\item 0.13\%, 1.27\% and 1.96\%.
\item Winsorising at 1\% still leaves large values in a distribution this heavy; ranking removes the scale entirely and
matches the saturating effect.
\item Cross-sectional ranks by date: robust to tails and outliers, comparable across dates.
\item 2.67\% held out; the truth 4.01\%.
\item Gain, permutation and drop-column rank $x_2$ first and $x_1$ fourth; SHAP ranks $x_2$, $x_1$, the copy, then $x_3$ and
$x_4$.
\item SHAP explains the fitted model, and the model uses the copy.
\item 5.4\% of the total (against 9.5\% for $x_1$): gain is measured on training data, where noise splits reduce the loss.
\item Drop-column importance (refit without the feature), and \omterm{def:ml:feature-engineering-selection-and-importance:mdi}{permutation importance} on held-out data only in part.
\item 0.76 and 0.41 points singly, 2.72 together.
\item 0.09 and $-0.02$ singly, 1.05 together.
\item $\{x_1, x_{1b}\}$ at a distance of 0.5 (correlation 0.95), all other features alone.
\item Dropped one copy, then the other, losing 1.05 of 2.67 points.
\item The lasso keeps $x_1$ (6.7 selected per half-sample; bound 2.3); boosting keeps $x_1$ and $x_2$ (bound 1.8); no noise
feature in either.
\item Their effect is only through the product: no linear screen sees it, and a greedy tree's first split on either
gains nothing.
\item \emph{Named result}: permutation and drop-column importance rank the signals $x_2, x_3, x_4, x_1$ (drop-column $x_2,
x_4, x_3, x_1$) and the copy far below; SHAP puts the copy third; the false-selection rate of \omterm{def:ml:feature-engineering-selection-and-importance:selection}{stability selection} at 0.6 is
zero for both selectors (bounds 2.3 and 1.8 false features), and both miss the interaction.
\item A portfolio manager: clustered drop-column importance (what the data need); a risk committee: SHAP for individual
decisions (what the model did).
\item Add candidate products or use a selector that can split on pairs (boosting with more trees, or interaction
constraints lifted), and give the pair a small main effect to find.
\item Importances should be computed on rolling windows; a feature's full-sample importance averages regimes.
\item Held-out clustered importances with the method named; the copies and their cluster; features that no method
ranks above noise; the selection procedure and its false-selection bound.
\item How much a fitted model, or the data, depend on a feature, in one sense that must be named.
\end{enumerate}
\end{solution}

\begin{solution}{iq:ml:feature-engineering-selection-and-importance:1}
Impurity importance sums the training-loss reductions of splits on a feature: biased toward features with many split
points and toward noise that was fitted. \omterm{def:ml:feature-engineering-selection-and-importance:mdi}{Permutation importance} measures the held-out score drop when the feature is
shuffled. Trust the held-out one, clustered by correlation.
\iqlookfor{training versus held-out, and the known biases of MDI.}
\end{solution}

\begin{solution}{iq:ml:feature-engineering-selection-and-importance:2}
Measure them as a group (permute or drop both): substitution hides each behind the other. If the group matters, keep one
or a combination (their mean or first principal component), chosen for stability and cost, not for its single
importance.
\iqlookfor{the \omterm{def:ml:feature-engineering-selection-and-importance:substitution}{substitution effect} and clustered importance.}
\end{solution}

\begin{solution}{iq:ml:feature-engineering-selection-and-importance:3}
The \omterm{def:ml:feature-engineering-selection-and-importance:shapley}{Shapley values} of the features for one prediction, with the model's expected output given subsets of features as
the game; they add up to the prediction minus the average. They mislead with correlated features (credit is shared by
a copy the model happens to use) and when read as causal effects rather than descriptions of the model.
\iqlookfor{the additivity property and the correlated-feature caveat.}
\end{solution}

\begin{solution}{iq:ml:feature-engineering-selection-and-importance:4}
By keeping only features selected in most half-samples; under exchangeability the expected number of noise features kept
is at most $q^2/((2\pi_{\mathrm{thr}} - 1)p)$. It does not control misses: features whose effect the base selector cannot
see (interactions, nonlinearities for a lasso) are never selected.
\iqlookfor{the bound and its blind spot.}
\end{solution}

\begin{solution}{iq:ml:feature-engineering-selection-and-importance:5}
Take logs, rank it by date or winsorise it at quantiles estimated on training data only, and check the transformed
feature's relation to the target is stable; consider a volume surprise (relative to its own history) rather than the
level.
\iqlookfor{robust transforms fitted without look-ahead.}
\end{solution}

\begin{solution}{iq:ml:feature-engineering-selection-and-importance:6}
Let $x_1, x_2$ be independent signs $\pm1$ and $y = x_1x_2$: each is uncorrelated with $y$, together they determine it.
Marginal screens (correlations, univariate tests, a lasso) discard both; screening must allow pairs, or a model that
can represent interactions must be the selector.
\iqlookfor{the XOR example and its consequence.}
\end{solution}
